\chapter{Modern Video Compression Standards}

Marco Cagnazzo - Multimedia Communications

\bigskip\hrule\bigskip

\section{The Hybrid Video Coding Paradigm}

\subsection{The Universal Hybrid Video Encoder}

From MPEG-1 to VVC, every standard uses the same fundamental \textbf{hybrid video coding} architecture.

\textbf{Core processes:}

\begin{itemize}
\item \textbf{Spatial prediction}, also called \textit{Intra prediction}.
\item \textbf{Temporal prediction}, also called \textit{Inter prediction}.
\item \textbf{Residual calculation}, obtained as the difference between the current signal and the prediction.
\item \textbf{Transform and quantization}, which decorrelate and compress the residual.
\item An \textbf{internal closed loop}, using inverse quantization and inverse transform, to reconstruct the same reference frames available at the decoder and prevent \textit{drift}.
\end{itemize}

\begin{calloutbox}{red!70!black}{Universal Hybrid Encoder}
\textbf{Statement:} Modern standards from MPEG-1 to VVC are based on spatial and temporal prediction, residual coding, transform, quantization, entropy coding, and an internal reconstruction loop.

\textbf{Interpretation:} The encoder must reconstruct its own decoded frames because future motion compensation uses decoded references, not original frames. This keeps the encoder and decoder synchronized and prevents drift.
\end{calloutbox}
\begin{center}\includegraphics[width=0.600\linewidth]{Pasted_image_20260506125516.png}\end{center}
\textit{Block diagram of the universal hybrid video encoder. The input frame $\hat{A}$ enters a summation node that produces residual $e_B$. The residual passes through two stages of $4 \times 4$ DCT, quantization, and CABAC/CAVLC to form the output bitstream. The reconstruction loop contains inverse quantization, inverse DCT, reconstructed residual $\hat{e}_B$, summation with prediction, frame buffer, deblocking filter, and reconstructed frames $\hat{f}_p$ / $\hat{f}_b$. A mode decision block selects between Intra Prediction and Motion Compensation, while Motion Estimation feeds Motion Compensation from the Frame Buffer.}
\subsection{Scope of the Standard}

If the fundamental architecture is universal, the key question is: what exactly is standardized?

\begin{center}\includegraphics[width=0.600\linewidth]{Pasted_image_20260506125634.png}\end{center}

\begin{calloutbox}{red!70!black}{Interoperability and Competition}
\textbf{Statement:} Standards exclusively define the bitstream syntax and the decoding process. Encoder architecture and optimization strategies remain open problems for industrial competition.

\textbf{Interpretation:} Any compliant decoder must decode a valid bitstream in the same way. Encoders, however, can compete through better mode decisions, faster search strategies, rate control, perceptual optimization, and hardware implementation.
\end{calloutbox}

\subsection{Codecs Historical Evolution}

To improve compression, standards continuously increase the encoder's \textbf{degrees of freedom}, especially partitioning options and prediction modes.

\begin{center}\includegraphics[width=0.600\linewidth]{Pasted_image_20260506125655.png}\end{center}

\begin{calloutbox}{red!70!black}{The Golden Rule}
\textbf{Statement:} Generational leaps typically halve the required bit-rate for identical perceptual quality, with an encoder computational complexity increase of approximately $10\times$.

\textbf{Interpretation:} Each new codec generation improves compression efficiency, but the encoder must search a much larger space of partitions, prediction modes, transforms, filters, and signaling choices.
\end{calloutbox}

\subsection{Market and Use Cases 2026 Scenarios}

Currently, multiple standards coexist. Codec adoption is driven by licensing models and hardware availability rather than strict technical superiority.

\begin{center}
\begin{tabular}{|l|l|l|}
\hline
\textbf{Ecosystem} & \textbf{Typical Codecs and Features} & \textbf{Market Share} \\ \hline
\textbf{Real-Time and Legacy} (video calls, old devices) & H.264/AVC: low-latency support and universal hardware compatibility. & H.264: approximately 80-90\% \\ \hline
\textbf{Broadcast and Mobile} (DVB-T2, premium 4K) & H.265/HEVC, H.264/AVC: high efficiency and mature hardware ecosystem. & HEVC: approximately 50-70\%; H.264: approximately 30-50\% \\ \hline
\textbf{Web Streaming} (YouTube, VOD) & AV1, H.264/AVC, H.265/HEVC: OTT-optimized, with AV1 as royalty-free. & H.264: approximately 60\%; AV1: approximately 20-30\%; HEVC: approximately 10-20\% \\ \hline
\textbf{Next-Gen} (8K, 360-degree video, VR) & H.266/VVC, AV1: maximum efficiency and high computational complexity. & AV1: growing; VVC: less than 5\% \\ \hline
\end{tabular}
\end{center}

\subsection{Video Codec Comparison}

The following comparisons use the UVG dataset and show the tradeoff between bitrate, objective quality, and encoding time.

\subsubsection{Test Sequence: HoneyBee}

\textbf{Test sequence:} HoneyBee, 1080p, 120 fps. \textbf{Target coding rate:} 2.0 Mbps.

\begin{center}
\begin{tabular}{|l|r|r|r|}
\hline
\textbf{Codec} & \textbf{Actual Rate} & \textbf{PSNR} & \textbf{Enc. Time (s)} \\ \hline
MPEG-2 & 7.06 Mbps & 36.58 dB & 2.18 \\ \hline
H.264 & 2.00 Mbps & 40.51 dB & 17.06 \\ \hline
HEVC & 1.98 Mbps & 41.69 dB & 109.45 \\ \hline
AV1 & 1.80 Mbps & 43.47 dB & 158.66 \\ \hline
VVC & 1.98 Mbps & 43.81 dB & 105.08 \\ \hline
\end{tabular}
\end{center}

\textit{A. Mercat, M. Viitanen, and J. Vanne, "UVG dataset: 50/120fps 4K sequences for video codec analysis and development," in Proc. ACM Multimedia Syst. Conf., Istanbul, Turkey, June 2020. Copyright: Digiturk 2013.}

\begin{calloutbox}{gray!60!black}{Figure 4 - HoneyBee visual comparison across codecs}
\textbf{Placeholder:} add image later from the original PDF. \textbf{Expected content:} HoneyBee sequence with purple flowers compared across MPEG-2, H.264, HEVC, AV1, and VVC. \textbf{Source cue:} \texttt{[Figura: HoneyBee sequence (purple flowers) compared across codecs.]}
\end{calloutbox}

\textbf{Visual comparison: HoneyBee}

\begin{itemize}
\item \textbf{MPEG-2:} Rate 7.06 Mbps, PSNR 36.58 dB. MPEG-2 is not able to attain the target bit-rate. Blocking artifacts are quite visible.
\item \textbf{H.264:} Rate 2.00 Mbps, PSNR 40.51 dB. H.264 attains the target bit-rate. Blocking artifacts are removed, but some artifacts are still visible.
\item \textbf{HEVC:} Rate 1.98 Mbps, PSNR 41.69 dB. Improved quality.
\item \textbf{AV1:} Rate 1.80 Mbps, PSNR 43.47 dB. Top quality.
\item \textbf{VVC:} Rate 1.98 Mbps, PSNR 43.81 dB. Top quality.
\end{itemize}

\subsubsection{Test Sequence: Bosphorus}

\textbf{Test sequence:} Bosphorus, 1080p, 120 fps. \textbf{VBV target:} 2.0 Mbps.

\begin{center}
\begin{tabular}{|l|r|r|r|l|}
\hline
\textbf{Codec} & \textbf{Actual Rate} & \textbf{PSNR} & \textbf{Enc. Time (s)} & \textbf{Playback} \\ \hline
MPEG-2 & 8.47 Mbps & 36.14 dB & 2.45 & [Play] \\ \hline
H.264 & 2.03 Mbps & 35.58 dB & 19.20 & [Play] \\ \hline
HEVC & 1.97 Mbps & 38.82 dB & 143.83 & [Play] \\ \hline
AV1 & 1.98 Mbps & 40.69 dB & 147.13 & [Play] \\ \hline
VVC & 1.96 Mbps & 41.00 dB & 139.08 & [Play] \\ \hline
\end{tabular}
\end{center}

\textit{A. Mercat, M. Viitanen, and J. Vanne, "UVG dataset: 50/120fps 4K sequences for video codec analysis and development," in Proc. ACM Multimedia Syst. Conf., Istanbul, Turkey, June 2020. Copyright: Digiturk 2013.}

\begin{calloutbox}{gray!60!black}{Figure 5 - Bosphorus visual comparison across codecs}
\textbf{Placeholder:} add image later from the original PDF. \textbf{Expected content:} Bosphorus sequence with a boat on water and a bridge, compared across MPEG-2, H.264, HEVC, AV1, and VVC. \textbf{Source cue:} \texttt{[Figura: Bosphorus sequence (boat on water with bridge) compared across codecs.]}
\end{calloutbox}

\textbf{Visual comparison: Bosphorus}

\begin{itemize}
\item \textbf{MPEG-2:} Rate 8.47 Mbps, PSNR 36.14 dB.
\item \textbf{H.264:} Rate 2.03 Mbps, PSNR 35.58 dB.
\item \textbf{HEVC:} Rate 1.97 Mbps, PSNR 38.82 dB.
\item \textbf{AV1:} Rate 1.98 Mbps, PSNR 40.69 dB.
\item \textbf{VVC:} Rate 1.96 Mbps, PSNR 41.00 dB.
\end{itemize}

\subsection{The Cost of Freedom}

The empirical data shows that newer standards achieve significantly better PSNR at the same bit-rate, but require ever increasing encoding time.

\begin{itemize}
\item \textbf{MPEG-2 underflow:} The encoder reaches its maximum quantization limit. It lacks the tools to discard further bits, causing severe target-rate overshoots of 7-8 Mbps.
\item \textbf{The RDO challenge:} VVC and AV1 have a massive tree of possible block partitions and prediction modes. Finding the optimal configuration under 2.0 Mbps while maximizing quality requires evaluating millions of combinations.
\end{itemize}

\textbf{How does the encoder make these decisions?}

\[
\rightarrow \text{Rate-Distortion Optimization}
\]

\section{Space Partitioning and RDO}

\subsection{Recall: Rate-Distortion Optimization}

From the previous lecture, the encoder minimizes the Lagrangian cost function for each block:

\[
J = D + \lambda R \tag{1}
\]

\begin{calloutbox}{red!70!black}{Rate-Distortion Optimization}
\textbf{Statement:} For each coding choice, the encoder minimizes

\[
J = D + \lambda R
\]

where $D$ is distortion, $R$ is rate, and $\lambda$ controls the tradeoff between quality and bitrate.

\textbf{Interpretation:} A mode is not selected only because it has low distortion or low rate. It is selected because it minimizes their weighted combination.
\end{calloutbox}

\textbf{The block partitioning problem:}

\begin{itemize}
\item Partitioning a block into sub-blocks is treated as a distinct coding mode.
\item The encoder evaluates the cost of the original block, $J_B$, against the sum of the sub-block costs:
\end{itemize}

\[
J_{split} = \sum J_{subblock}
\]

\begin{itemize}
\item If $J_{split} < J_B$, the split is accepted.
\item \textit{The evolution of video standards is largely the evolution of allowable split geometries.}
\end{itemize}

\subsection{H.264/AVC: The Macroblock Era}

H.264 relies on a fixed and rigid structure, which limits the RDO search space.

\begin{itemize}
\item \textbf{Base unit:} the $16 \times 16$ \textbf{Macroblock} (MB).
\item \textbf{Variable Block-Size (VBS):} the MB can be partitioned into $16 \times 8$, $8 \times 16$, or $8 \times 8$ blocks to better capture motion boundaries.
\item \textbf{Sub-partitions:} $8 \times 8$ blocks can be further divided down to $4 \times 4$.
\end{itemize}

\begin{calloutbox}{red!70!black}{Limitation of the Macroblock}
\textbf{Statement:} A maximum block size of $16 \times 16$ is highly inefficient for High Definition (1080p) or 4K video, where smooth background areas such as sky or walls can span hundreds of pixels.

\textbf{Interpretation:} Large smooth regions should be represented by large blocks. If the standard only permits small blocks, the encoder must spend unnecessary bits on partition signaling and local prediction.
\end{calloutbox}

\subsection{H.265/HEVC: The Quad-Tree Paradigm}

HEVC abandons macroblocks in favor of the \textbf{Coding Tree Unit (CTU)} to handle UHD resolutions.

\begin{itemize}
\item \textbf{Larger base size:} CTUs can be up to $64 \times 64$ pixels.
\item \textbf{Quad-Tree split:} a CTU is recursively split into four square \textbf{Coding Units} (CUs), down to $4 \times 4$ CUs.
\item \textbf{Functional separation:} a CU can be symmetrically or asymmetrically divided into \textbf{Prediction Units} (PUs) and \textbf{Transform Units} (TUs).
\end{itemize}

\begin{center}\includegraphics[width=0.600\linewidth]{Pasted_image_20260506125829.png}\end{center}
\begin{center}\includegraphics[width=0.600\linewidth]{Pasted_image_20260506125913.png}\end{center}
\begin{center}\includegraphics[width=0.600\linewidth]{Pasted_image_20260506125928.png}\end{center}
\subsection{H.266/VVC: QT + Multi-Type Tree (MTT)}

VVC introduces extreme flexibility and eliminates the strict separation between CUs, PUs, and TUs for most modes.

\begin{itemize}
\item \textbf{Massive base size:} CTUs can be up to $128 \times 128$ pixels.
\item \textbf{QT + MTT:} Quad-Tree leaves can be further partitioned using a \textbf{Multi-Type Tree}.
\item \textbf{New splits:} includes \textit{Binary} splits, which are symmetric halves, and \textit{Ternary} splits, with a $1:2:1$ ratio.
\item Rectangular blocks become fundamental for encoding object edges accurately.
\end{itemize}
\begin{center}\includegraphics[width=0.600\linewidth]{Pasted_image_20260506125946.png}\end{center}
\subsection{AV1: 10-Way Superblock Partitioning}

AV1 uses a similar highly flexible hierarchical approach.

\begin{itemize}
\item \textbf{Base unit:} the \textit{Superblock}, sized up to $128 \times 128$.
\item \textbf{10-way split:} a block can be partitioned using 10 different geometrical patterns.
\item The available geometries include $4:1$ and $1:4$ rectangular splits, creating highly elongated blocks for sharp directional textures.
\item Square blocks can be recursively partitioned.
\end{itemize}

\begin{center}\includegraphics[width=0.600\linewidth]{Pasted_image_20260506130003.png}\end{center}

\subsection{The Algorithmic Burden of RDO}

\begin{calloutbox}{red!70!black}{Decision Tree Search}
\textbf{Statement:} To find the optimal partition, the encoder explores the block geometry using a \textbf{Depth-First Search (DFS)} algorithm while optimizing the cost

\[
J = D + \lambda R.
\]

\textbf{Interpretation:} More partition freedom improves compression only if the encoder can afford to search the corresponding decision tree.
\end{calloutbox}

Idea: use the encoded block to predict the next block -\textgreater{} sequential The evaluation is a massive sequential process:

\begin{enumerate}
\item \textbf{Top-down exploration:} try coding the block as a whole, then try splitting it.
\item \textbf{Causal dependency:} sub-blocks are tested sequentially in Z-scan order. A block cannot be predicted until its top and left neighboring blocks have been fully \textit{reconstructed}, so their pixels are frozen.
\item \textbf{Inner loop for prediction:} for each block, exhaustively test Intra and Inter modes to find the local minimum $J_{CU}$.
\item \textbf{Bottom-up decision:} compare the total cost. If
\end{enumerate}

\[
\sum J_{children} < J_{parent},
\]

the split is kept.

\textbf{Conclusion:} The theoretical state space of a CTU is enormous, but reference models traverse it sequentially. Commercial encoders drastically reduce this process using \textbf{fast heuristics} and \textbf{machine learning} to prune the tree early.

\subsection{Pruning the Decision Tree: Fast Heuristics}

To mitigate the extreme complexity of RDO, production encoders such as VVenC and SVT-AV1 implement aggressive early termination strategies.

\textbf{Common pruning approaches:}

\begin{itemize}
\item \textbf{Spatial/texture analysis:} evaluate local variance and directional gradients. If a $64 \times 64$ Coding Unit has low variance, as in a flat background, deeper QT/MTT partitions are immediately bypassed.
\item \textbf{Temporal correlation:} constrain the RDO search space by exploiting the partition depth of the co-located CU in the reference frame.
\item \textbf{Machine learning classifiers:} use lightweight CNNs or gradient boosting decision trees, such as LightGBM, to predict split probabilities directly from raw pixel features and discard highly improbable geometries before any RDO calculation.
\end{itemize}

\section{Spatial and Temporal Prediction}

\subsection{Intra Prediction}

\textbf{Intra prediction} exploits spatial redundancy by extrapolating pixels from already-decoded neighboring blocks.
\begin{center}\includegraphics[width=0.806\linewidth]{Pasted_image_20260506130740.png}\end{center}

\textbf{The evolution of freedom:}

\begin{itemize}
\item \textbf{H.264/AVC:} 4 prediction types for $16 \times 16$ blocks, and 9 types for $4 \times 4$ blocks: 8 directions plus DC.
\item \textbf{H.265/HEVC:} expands to 35 modes: 33 directions plus DC and Planar, to handle complex textures.
\item \textbf{H.266/VVC:} expands to 65 directional modes, plus wide-angle modes for non-square blocks.
\end{itemize}

The number of possible predictors grows exponentially with the number of signaling bits:

\[
P \propto 2^B.
\]

\begin{center}\includegraphics[width=0.806\linewidth]{Pasted_image_20260506130804.png}\end{center}
\textit{Horizontal predictor}

\begin{center}\includegraphics[width=0.806\linewidth]{Pasted_image_20260506130818.png}\end{center}
\textit{Vertical predictor}

\begin{center}\includegraphics[width=0.806\linewidth]{Pasted_image_20260506130833.png}\end{center}
**

\begin{center}\includegraphics[width=0.806\linewidth]{Pasted_image_20260506130853.png}\end{center}
\begin{center}\includegraphics[width=0.806\linewidth]{Pasted_image_20260506130918.png}\end{center}
\textit{we can split it in 16 4x4 blocks}

\begin{center}\includegraphics[width=0.600\linewidth]{Pasted_image_20260506131234.png}\end{center}
\begin{center}\includegraphics[width=0.600\linewidth]{Pasted_image_20260506131249.png}\end{center}
\begin{center}\includegraphics[width=0.600\linewidth]{Pasted_image_20260506131301.png}\end{center}
\begin{center}\includegraphics[width=0.600\linewidth]{Pasted_image_20260506131310.png}\end{center}

For each sub block we have 9 options. We can build 9 times 16 possible combinations. We have $9^{16}$ possible predictions ideally. We do it sequentially to have fewer combination.
\subsection{The Cost of Freedom: Predicting the Predictor}

\begin{calloutbox}{red!70!black}{Overhead Problem}
\textbf{Statement:} If VVC has 67 Intra modes, sending the mode index requires approximately 7 bits per block. For tiny $4 \times 4$ blocks, this signaling overhead would completely destroy the compression gains.

\textbf{Interpretation:} Increasing the number of modes improves prediction quality, but the selected mode must also be transmitted. The mode index itself can become too expensive.
\end{calloutbox}

\textbf{Solution: Most Probable Mode (MPM)}

\begin{itemize}
\item The encoder predicts the mode itself using the spatial context.
\item A list of \textbf{Most Probable Modes} (MPM) is built using the directions chosen by the top and left neighbors.
\item HEVC uses a 3-candidate MPM list.
\item VVC expands this to a 6-candidate MPM list.
\item If the chosen mode is in the MPM list, it is coded with a very short index, for example 1 or 2 bits.
\end{itemize}

\subsection{Inter Prediction: Fractional Motion Compensation}

\textbf{Inter prediction} removes temporal redundancy through \textbf{Motion Estimation} and \textbf{Motion Compensation} (ME/MC).

\begin{itemize}
\item Objects rarely move by exact integer pixel amounts between frames.
\item \textbf{H.264/HEVC/VVC:} use quarter-pel, or $1/4$-pixel, fractional precision. The encoder generates sub-pixels using interpolation filters, such as the 6-tap filters in H.264.
\item \textbf{AV1:} pushes precision further to $1/8$-pixel accuracy and uses dynamically selectable interpolation filters, such as sharp and smooth filters.
\end{itemize}

\subsection{Advanced Motion Compensation}

Modern standards generalize temporal distance. A frame is no longer restricted to predicting from its immediate chronological neighbor.

\begin{itemize}
\item \textbf{Multiple reference frames:} slices can maintain one P-list or two B-lists of reference images kept in the \textbf{Decoded Picture Buffer} (DPB).
\item \textbf{Generalized B-slices:} the encoder can flexibly select predictors from these lists and combine them using weighted averages.
\end{itemize}

\begin{center}\includegraphics[width=0.600\linewidth]{Pasted_image_20260506133134.png}\end{center}

\subsection{Vector Prediction: From Skip to Merge}

Just like Intra modes, \textbf{Motion Vectors} (MVs) must be predicted to save bits.

\begin{itemize}
\item \textbf{H.264 Skip/Direct:} uses a simple median of spatial neighbors. If the median is correct, 0 bits are spent on MVs.
\item \textbf{HEVC/VVC Merge Mode:} replaces the median with a dynamic \textit{candidate list}.
\item The list aggregates 5 spatial candidates from adjacent blocks and 1 temporal candidate from the co-located block in the reference frame.
\item The encoder transmits only the index of the best candidate in the list, sharing both the MV and the reference frame index.
\end{itemize}

\subsection{Beyond Translation: Modern Inter Tools}

To squeeze the final bits of temporal redundancy out of the signal, AV1 and VVC abandon the assumption that blocks only move rigidly in 2D space.

\textbf{Affine Motion Prediction:}

\begin{itemize}
\item Allows blocks to rotate, zoom, and shear.
\item Uses 4-parameter or 6-parameter models instead of a simple $(x, y)$ vector.
\end{itemize}

\textbf{Compound Prediction:}

\begin{itemize}
\item AV1 allows Inter and Intra predictions to be combined within the same block.
\end{itemize}

\begin{center}\includegraphics[width=0.600\linewidth]{Pasted_image_20260506134619.png}\end{center}

\section{Transforms, Quantization, and Filters}

\subsection{Residual Coding: The Evolution of Transforms}

After prediction, the residual error must be decorrelated into the frequency domain.

\textbf{The push for larger and flexible transforms:}

\begin{itemize}
\item \textbf{H.264/AVC:} relies on a $4 \times 4$ integer-coefficient approximation of the DCT, perfectly synchronized to avoid precision errors between encoder and decoder.
\item \textbf{H.265/HEVC:} expands transforms up to $32 \times 32$ to handle 4K resolutions. It also introduces a DST, or \textit{Discrete Sine Transform}, approximation specifically for $4 \times 4$ Intra blocks, improving compression by approximately $1\%$.
\item \textbf{H.266/VVC:} introduces \textbf{Multiple Transform Selection} (MTS) and non-square transforms to match the rectangular partitions created by QT+MTT.
\end{itemize}

\subsection{Quantization and Entropy Coding}

\textbf{Quantization} is the only lossy step in the hybrid encoder.

\textbf{Quantization Parameter (QP):}

\begin{itemize}
\item In H.264/HEVC, the quantization step size doubles for every increment of 6 in the \textbf{Quantization Parameter} (QP).
\item Empirically, a unitary QP increment increases the rate by approximately $12.5\%$.
\end{itemize}

\begin{calloutbox}{red!70!black}{Quantization Parameter Behavior}
\textbf{Statement:} In H.264/HEVC, the quantization step size doubles whenever QP increases by 6. Empirically, increasing QP by 1 changes the rate by approximately $12.5\%$.

\textbf{Interpretation:} QP is the main knob for rate control. Higher QP means coarser quantization, lower bitrate, and higher distortion.
\end{calloutbox}

\textbf{Lossless entropy coding:}

\begin{itemize}
\item Modern standards are standardized around \textbf{CABAC}, or \textit{Context-Adaptive Binary Arithmetic Coding}.
\item CABAC transforms syntax elements into binary formats and dynamically adapts probabilities based on previously decoded spatial contexts.
\item CABAC yields a 5\%-15\% rate reduction over older variable-length codes such as CAVLC.
\end{itemize}

\subsection{The Need for In-Loop Filtering}

Because transform and quantization are applied independently to disjoint blocks, strong discontinuities, called \textbf{blocking artifacts}, appear at block boundaries at low bit-rates.

\begin{center}\includegraphics[width=0.600\linewidth]{Pasted_image_20260506134648.png}\end{center}

If left uncorrected, these artifacts are saved in the Decoded Picture Buffer (DPB) and permanently damage future motion compensations. But here we know where we have those artifacts: around the contours.

\subsection{Evolution of In-Loop Filters}

To fix quantization damage, standards mandate complex \textbf{normative filters} inside the reconstruction loop.

\begin{itemize}
\item \textbf{H.264/AVC Deblocking:} analyzes edges between $4 \times 4$ blocks. Filtering strength adapts dynamically to the coding mode, motion vectors, and quantization step of neighboring blocks.
\item \textbf{H.265/HEVC Deblocking + SAO:} applies deblocking only on an $8 \times 8$ grid to reduce complexity. It introduces \textbf{Sample Adaptive Offset} (SAO), which categorizes pixels by edges and bands, then adds offsets to correct ringing artifacts.
\end{itemize}

\subsection{VVC and AV1 Filters}

As transform sizes grew and bit-rates decreased, deblocking and SAO were no longer sufficient. Modern codecs dedicate major computational effort to pixel restoration.

\begin{itemize}
\item \textbf{H.266/VVC:} adds the \textbf{Adaptive Loop Filter} (ALF) and \textbf{Luma Mapping with Chroma Scaling} (LMCS), specifically designed to preserve details in HDR, or \textit{High Dynamic Range}, video.
\item \textbf{AV1:} implements a strong post-processing pipeline:
\begin{itemize}
\item \textbf{Directional Filters (CDEF):} analyze gradient directions to smooth edges without blurring textures.
\item \textbf{Loop Restoration Filters:} Wiener filters applied over large windows.
\item \textbf{Film Grain Synthesis:} denoises the image before encoding, then synthesizes artificial grain at the decoder to recreate the cinematic look at zero bitrate cost.
\end{itemize}
\end{itemize}

\section{Network Abstraction and Parallelism}

\subsection{Transmission and Network Errors}

\begin{calloutbox}{red!70!black}{Fragility of Compression}
\textbf{Statement:} A 4K video compressed into a perfect 2 Mbps mathematical stream can be severely corrupted if only a few bits are flipped.

\textbf{Interpretation:} Entropy-coded streams are compact but fragile. A bit error can desynchronize CABAC contexts and then propagate through spatial prediction and motion compensation.
\end{calloutbox}

\textbf{Different "pipes", same vulnerability:}

\begin{itemize}
\item \textbf{DVB-T broadcast:} a continuous stream subject to electromagnetic noise.
\item \textbf{Internet, using RTP/UDP or TCP/DASH:} a stream cut into packets with MTU around 1500 bytes; packets can arrive late or be lost.
\item \textbf{Disk storage, such as MP4:} requires specific indexes to seek.
\end{itemize}

\textbf{Consequence:} CABAC loses context. Errors propagate spatially and temporally through Motion Compensation, creating a "green and gray mosaic".

\subsection{Pictures and Slices}

To limit error propagation, pictures can be divided into \textbf{slices}.

At each new slice, Intra prediction is stopped and CABAC is reset, so parallel encoding and decoding of slices becomes possible.

\begin{center}\includegraphics[width=0.600\linewidth]{Pasted_image_20260506135622.png}\end{center}

\subsection{The Solution: Separation of Concerns}

Starting from H.264, video standards separate \textbf{compression} from \textbf{transport}.

\textbf{The factory analogy:}

\begin{center}\includegraphics[width=0.600\linewidth]{Pasted_image_20260506135653.png}\end{center}

\begin{itemize}
\item \textbf{Video Coding Layer (VCL):} the production department. It handles pure compression: transforms, prediction, and RDO. It generates the \textit{slices} of each compressed picture.
\item \textbf{Network Abstraction Layer (NAL):} the shipping department. It takes VCL slices, encapsulates them, and adds metadata.
\item \textbf{Goal:} provide a standard interface to map video over RTP/IP for real-time communication, ISO MP4 for storage, or MPEG-2 TS for diffusion.
\end{itemize}

\subsection{NAL Units}

The VCL encodes each picture into one or more slices. The NAL encapsulates the slices into \textbf{VCL NALUs}, one per slice, and generates \textbf{Non-VCL NALUs}, such as SPS, PPS, and SEI.

\begin{itemize}
\item \textbf{VCL NALU:} contains modes, MVs, and quantized coefficients of a slice. It can be \textbf{Instantaneous Decoding Refresh}, meaning Intra with CABAC reset, or non-IDR.
\item \textbf{SPS:} \textbf{Sequence Parameter Set}, a non-VCL NALU containing important encoding parameters such as picture size and color depth.
\item \textbf{PPS:} \textbf{Picture Parameter Set}, containing per-picture coding options specific to each codec.
\item \textbf{SEI:} \textbf{Supplemental Enhancement Information}, not necessary for decoding but useful for specific tasks such as subtitles, HDR information, user data, and similar metadata.
\item \textbf{AUD:} \textbf{Access Unit Delimiter}, an optional NALU that signals a new picture.
\end{itemize}

\subsection{NALUs and Access Units}

\begin{calloutbox}{red!70!black}{Access Unit}
\textbf{Statement:} An \textbf{Access Unit} (AU) is a logical grouping of all NALUs required to decode a single frame. The decoder identifies its boundaries implicitly.

\textbf{Interpretation:} The AU is the frame-level unit of decoding, while NALUs are the transport and syntax units that compose it.
\end{calloutbox}

How is this sequence of NALUs delivered?

\begin{itemize}
\item If the transport medium is a continuous "blind" tube, such as MPEG-TS or raw \texttt{.264} files, use the \textbf{Byte-Stream format}.
\item If the transport layer already provides boundaries, such as an MP4 container, use the \textbf{Packetized format}.
\end{itemize}

\begin{center}\includegraphics[width=0.600\linewidth]{Pasted_image_20260506140326.png}\end{center}

\subsection{Format 1: Annex-B Byte-Stream}

\textbf{Mechanism:} a Start Code, for example \texttt{0x00000001}, is injected before every NALU.

\textbf{Emulation prevention:} fake bytes are injected into the payload to prevent valid video data from accidentally mimicking a Start Code.

\textbf{Pro:} Annex-B is well suited to continuous streaming or broadcasting, where packet boundaries do not exist natively.

\textbf{Con:} it introduces CPU overhead for emulation prevention and re-synchronization. The decoder must scan every byte sequentially to find the next Start Code.

\begin{center}\includegraphics[width=0.600\linewidth]{Pasted_image_20260506140340.png}\end{center}

\subsection{Format 2: Packetized Length-Prefixed}

\textbf{Mechanism:} a \textbf{Length Field}, usually 4 bytes, is written exactly before the NALU header.

\textbf{No Start Codes:} emulation prevention is not strictly required for boundary detection.

\textbf{Pro:} this format is highly software-friendly. The parser reads the length and performs an $O(1)$ memory jump directly to the next NALU.

\textbf{Con:} it is useless if the stream is corrupted or the external container is damaged, because synchronization is lost.

\begin{center}\includegraphics[width=0.600\linewidth]{Pasted_image_20260506140352.png}\end{center}

\subsection{The Evolution: HEVC, VVC, and AV1}

\textbf{H.265/HEVC and H.266/VVC:}

\begin{itemize}
\item Maintain the Annex-B / Packetized dualism.
\item Expand the NALU header from 1 byte to \textbf{2 bytes} to include \textit{Temporal ID} and \textit{Layer ID} for spatial and temporal scalability.
\end{itemize}

\textbf{AV1: breaking the paradigm}

\begin{itemize}
\item AV1 abandons the MPEG legacy. NALUs are called \textbf{OBUs}, or \textit{Open Bitstream Units}.
\item \textbf{No Start Codes:} AV1 completely drops Annex-B and emulation prevention to save CPU cycles.
\item AV1 exclusively uses a length-prefixed approach, even for raw bitstreams.
\end{itemize}

\subsection{Network Friendliness in Action}

The NAL enables intelligent network operations without decoding. The NAL header and structure allow intermediate nodes, called \textbf{Media Aware Network Elements} (MANEs), to understand and manage the video stream.

\textbf{1. Router inspection (MANE):}

\begin{itemize}
\item MANEs can read the 1-byte or 2-byte header; this is a form of \textit{Deep Packet Inspection}.
\item \textbf{Smart dropping:} if the network is congested, the router can selectively drop Non-Reference B-frames without touching critical IDR frames or parameter sets.
\item Protection, such as CRC, can be added to sensitive data such as SPS.
\end{itemize}

\textbf{2. RTP transport and fragmentation:}

\begin{itemize}
\item If a VCL NALU, such as an entire I-frame slice, exceeds the network MTU of 1500 bytes, it cannot be transmitted as-is.
\item The RTP layer slices large NALUs into \textbf{Fragmentation Units} (FUs) and reassembles them at the receiver, keeping the video decoder unaware of the network constraints.
\end{itemize}

\begin{center}\includegraphics[width=0.600\linewidth]{Pasted_image_20260506140426.png}\end{center}
\begin{center}\includegraphics[width=0.600\linewidth]{Pasted_image_20260506140532.png}\end{center}
\begin{center}\includegraphics[width=0.600\linewidth]{Pasted_image_20260506140606.png}\end{center}
\subsection{Resynchronization: Stream Structure}

\begin{center}\includegraphics[width=0.600\linewidth]{Pasted_image_20260506140635.png}\end{center}

Decoding may restart at a \textbf{Random Access Point} (RAP).

\textbf{IDR (Instantaneous Decoding Refresh):}

An IDR is a strict barrier that completely clears the Decoded Picture Buffer (DPB). No pictures following the IDR in decoding order can reference pictures preceding it.

\begin{itemize}
\item \textbf{Pros:} guarantees a clean restart. It is ideal for the absolute beginning of a video sequence or hard scene cuts.
\item \textbf{Cons:} breaks temporal prediction across the barrier, reducing compression efficiency. It cannot exploit leading pictures for better referencing.
\end{itemize}

\textbf{CRA (Clean Random Access):}

CRA was introduced in HEVC. It is an Intra picture that allows \textbf{RASL} (\textit{Random Access Skipped Leading}) frames to reference pictures located before the CRA.

\begin{itemize}
\item \textbf{Scenario 1, normal playback:} the decoder has past frames in memory. RASL pictures are fully decoded, maximizing compression efficiency.
\item \textbf{Scenario 2, seeking:} decoding starts directly at the CRA. Past references are missing, so the decoder safely discards, or skips, the RASL pictures and decodes the rest cleanly.
\item \textbf{Comparison:} H.264 lacked this feature, forcing frequent use of costly IDR frames.
\end{itemize}

\textbf{BLA (Broken Link Access):}

BLA functions similarly to CRA, but explicitly signals that the required reference pictures for its leading frames are destroyed.

\begin{itemize}
\item \textbf{Origin:} BLA is typically generated by a \textit{secondary encoder} in network nodes, such as transcoders, packagers, and ad-insertion splicers, rather than by the original content encoder.
\item \textbf{Mechanism:} when a stream is cut and spliced at a CRA, the network node alters the NALU header from CRA to BLA. The decoder reads this and immediately drops the broken RASL pictures without attempting reference lookups.
\end{itemize}

\subsection{Hardware Friendliness: Tiles}

\begin{itemize}
\item The CTU may be grouped into rectangular regions called \textbf{tiles}.
\item Tiles are self-contained and independent from slices.
\item A slice may contain several tiles, or a tile may contain several slices.
\item One of the main targets of tiles is to allow \textbf{parallel encoding and decoding} across multi-core processors.
\item Another application is viewport priority in 360-degree video.
\end{itemize}

\begin{center}\includegraphics[width=0.600\linewidth]{Pasted_image_20260506140722.png}\end{center}

\subsection{Tiles in 360-Degree Video}

Tiles are the foundational technology for \textbf{Viewport-Adaptive Streaming} in Virtual Reality, especially MPEG-OMAF.

\begin{center}\includegraphics[width=0.600\linewidth]{Pasted_image_20260506140746.png}\end{center}

\begin{itemize}
\item \textbf{The problem:} a VR headset only displays approximately 15\% of the full 360-degree sphere, the Field of View. Transmitting the unseen 85\% at 8K is a massive waste of bandwidth.
\item \textbf{The tile solution:} the sphere is projected and encoded as a grid of independent HEVC tiles.
\item \textbf{Dynamic stitching:} the client requests only the tiles inside the current viewport at high resolution, and requests the background tiles at very low resolution as a fallback.
\item \textbf{Why tiles?} Motion vectors cannot cross tile boundaries, so tiles are independent and can be mixed and matched at different qualities.
\end{itemize}

\begin{calloutbox}{red!70!black}{Bandwidth Savings With Tiles}
\textbf{Statement:} Tile-based adaptive streaming can reduce bandwidth consumption by 50\%-60\% without degrading the user's perceived quality.

\textbf{Interpretation:} The system spends high bitrate only where the user is looking and sends low-quality fallback content outside the viewport.
\end{calloutbox}

\subsection{Wavefront Parallel Processing (WPP)}

\begin{itemize}
\item A finer degree of parallelism can be achieved using \textbf{Wavefront Parallel Processing} (WPP).
\item Slices are divided into rows of CTUs.
\item The encoding of a row may start as soon as two CTUs of the previous row have been encoded.
\item There is a guaranteed delay of 2 CTUs between successive rows, preserving spatial prediction contexts.
\end{itemize}

![[Pasted image 20260506140812.png]]`

\subsection{Summary of the Hybrid Architecture}

\textbf{The cost of compression:}

\begin{itemize}
\item Standards moved from fixed $16 \times 16$ blocks to extreme QT+MTT geometries up to $128 \times 128$.
\item Intra directions expanded from 9 to 67, and sub-pixel affine motion was introduced.
\item Complex in-loop filters, such as ALF and Restoration filters, were deployed to mask quantization artifacts.
\item NALUs, mappings, tiles, and WPP were engineered to make the compressed video transportable and parallelizable.
\end{itemize}

\begin{calloutbox}{red!70!black}{The Next Frontier}
\textbf{Statement:} The hybrid paradigm has arguably reached its mathematical asymptote with VVC. The next major leap requires abandoning discrete mathematical transforms in favor of end-to-end differentiable neural architectures.

\[
\rightarrow \textit{Next Lectures: Neural Video Coding}
\]

\textbf{Interpretation:} Further compression gains may require learned representations and neural codecs rather than simply adding more handcrafted prediction, transform, partitioning, and filtering tools.
\end{calloutbox}

\section{Final Summary Table}

\begin{center}
\begin{tabular}{|l|l|l|l|}
\hline
\textbf{Tool / Standard} & \textbf{Key idea} & \textbf{Main parameters} & \textbf{Typical use} \\ \hline
Hybrid video encoder & Prediction plus residual transform coding inside a closed reconstruction loop. & Intra/Inter prediction, residual, transform, quantization, entropy coding, inverse loop. & Universal architecture from MPEG-1 to VVC. \\ \hline
Scope of standard & Standardize bitstream syntax and decoding, not encoder optimization. & Decoder behavior, syntax elements, conformance rules. & Interoperability and industrial competition. \\ \hline
Rate-Distortion Optimization & Select the mode minimizing $J = D + \lambda R$. & Distortion $D$, rate $R$, multiplier $\lambda$. & Partitioning, prediction mode choice, transform and quantization decisions. \\ \hline
H.264/AVC macroblocks & Fixed $16 \times 16$ MB with variable block-size partitions. & $16 \times 16$, $16 \times 8$, $8 \times 16$, $8 \times 8$, down to $4 \times 4$. & Legacy and real-time video with broad hardware support. \\ \hline
H.265/HEVC CTU & Quad-tree Coding Tree Unit for UHD content. & CTU up to $64 \times 64$, CUs, PUs, TUs. & Efficient HD, 4K, broadcast, and mobile compression. \\ \hline
H.266/VVC QT+MTT & Quad-tree plus binary and ternary multi-type splits. & CTU up to $128 \times 128$, binary split, ternary $1:2:1$ split. & Maximum coding efficiency with high encoder complexity. \\ \hline
AV1 superblock partitioning & Flexible 10-way partitioning of $128 \times 128$ superblocks. & Square, rectangular, $4:1$, and $1:4$ split options. & Web streaming and royalty-free distribution. \\ \hline
Intra prediction & Predict current block from already-decoded spatial neighbors. & H.264 4/9 modes, HEVC 35 modes, VVC 65+ modes. & Spatial redundancy removal. \\ \hline
Most Probable Mode & Predict the Intra mode using top and left neighboring modes. & HEVC 3-candidate list, VVC 6-candidate list. & Reduce signaling overhead for many Intra modes. \\ \hline
Fractional ME/MC & Estimate motion at sub-pixel precision. & Quarter-pel for H.264/HEVC/VVC, eighth-pel for AV1. & Temporal redundancy removal. \\ \hline
Merge mode & Transmit an index into a motion-vector candidate list. & Spatial candidates, temporal co-located candidate, reference frame index. & MV compression in HEVC/VVC. \\ \hline
Affine motion & Model rotation, zoom, and shear instead of simple translation. & 4-parameter or 6-parameter affine models. & Advanced Inter prediction in modern codecs. \\ \hline
Transform coding & Decorrelate prediction residuals in the frequency domain. & H.264 $4 \times 4$ integer DCT, HEVC up to $32 \times 32$, VVC MTS. & Compact residual representation. \\ \hline
Quantization & Lossy reduction of transform coefficient precision. & QP, step size doubling every +6 QP, approximately 12.5\% rate change per QP. & Main bitrate and distortion control. \\ \hline
CABAC & Context-adaptive binary arithmetic coding. & Binary syntax elements, adaptive spatial contexts. & Lossless entropy coding with 5\%-15\% rate savings over older VLC. \\ \hline
In-loop filters & Normative filtering inside the reconstruction loop. & Deblocking, SAO, ALF, LMCS, CDEF, restoration filters. & Prevent artifacts from entering reference frames. \\ \hline
VCL and NAL & Separate compression from transport encapsulation. & VCL slices, NAL units, SPS, PPS, SEI, AUD. & Mapping video to RTP/IP, MP4, MPEG-2 TS, and other systems. \\ \hline
Annex-B & Byte-stream format with Start Codes. & \texttt{0x00000001}, emulation prevention. & Continuous streams and broadcasting. \\ \hline
Packetized format & Length-prefixed NALU format. & Usually 4-byte length field. & MP4-like containers and software-friendly parsing. \\ \hline
AV1 OBU & AV1 replaces NALUs with Open Bitstream Units. & Length-prefixed OBUs, no Annex-B Start Codes. & AV1 raw bitstreams and web-oriented delivery. \\ \hline
Random access & Restart decoding at well-defined access points. & IDR, CRA, RASL, BLA. & Seeking, splicing, scene cuts, clean recovery. \\ \hline
Tiles & Self-contained rectangular CTU regions. & Tile grid, independent motion constraints. & Parallel processing and viewport-adaptive 360-degree streaming. \\ \hline
WPP & Row-level CTU parallelism with a two-CTU delay. & CTU rows, diagonal wavefront, 2-CTU offset. & Multi-core encoding and decoding. \\ \hline
\end{tabular}
\end{center}

